% page 54 of the facsimile (image p-053 of the scan)
% block 162.6 x 237.7 mm ; left margin 8.0 mm
\pgc{-12.700mm}{0.000mm}{
\l{\cel{6}46}
\l{}
\l{}
\l{\sou{astro}. (\sou{astron.}) Ciel-korpo. - DEFIS.}
\l{}
\l{\sou{astriktar}.\cel{2}(\sou{trans.}) Densigar ula tisui organika. (Aluno esas}
\l{\cel{3}astriktiva). - DEFIS.}
\l{}
\l{\sou{astrocentro}. (\sou{biol}.) Fuzelo akromata di la nukleo. - DEFIRS.}
\l{}
\l{\sou{astrofiziko}. Studio di la astri per la metodi diversa di fiziko}
\l{\cel{3}- DEFIRS.}
\l{}
\l{\sou{astroida}.\cel{2}Astro-forma, stelo-forma. - DEF.}
\l{}
\l{\sou{astrolabo}.\cel{2}Instrumento uzata por mezurar la alteso-grado di}
\l{\cel{3}la astri super horizonto. - DEFIRS.}
\l{\sou{astrologo}. La persono qua praktikas astrologio. - DEFIRS.}
\l{\sou{astrologio}. Exploro di la astri pri lia influo asertita di la}
\l{\cel{3}homo-fato. - DEFIRS.}
\l{}
\l{\sou{astronomo}. La persono qua konsakras su ad astronomio. - DEFIRS.}
\l{}
\l{\sou{astronomio}.\cel{2}Cienco pri la astri. - DEFIRS.}
\l{}
\l{\sou{astrosfero}.\cel{2}(\sou{biol}.) Sfero atraktiva. -DEFIRS.}
\l{}
\l{\sou{asumar}.\cel{2}(\sou{trans.})\cel{2}Aceptar responsor pri (ulo). - EFIS.}
\l{}
\l{\sou{asunciono}. (\sou{religio katolika}). Festo-dio institucita da la Ekle-}
\l{\cel{3}zio katolika, por honorizar la transporto mirakla di la kor-}
\l{\cel{3}po e di la anmo di la Virgino aden cielo. - EFIS.}
\l{}
\l{\sou{- at -}\cel{1}. Sufixo qua indikas la participo prezenta pasiva.}
\l{}
\l{\sou{atakar}.\cel{2}(\sou{trans., per, ye)}\cel{2}Komencar la frapo (di enemiko, di}
\l{\cel{3}adverso). - DEFIRS.}
\l{}
\l{\sou{ataraxio}. Quieteso di la psiko quan trublas nula deziro, nula}
\l{\cel{3}timo. - DEFIRS.}
\l{}
\l{\sou{atasheo}.- (\sou{diplomaco)}.\cel{2}La persono qua, en ambasadeyo, esas,}
\l{\cel{3}sur la skalo hierarkiala, nemediate dop la sekretario. -}
\l{\cel{3}DEFIR.}
\l{}
\l{\sou{atavismo}.\cel{2}Herediveso di ula karakteri korpala o mentala qui,}
\l{\cel{3}en ula kazi, devenas de plura generacioni. - DEFIRS.}
\l{}
\l{\sou{ataxio}. (\sou{patol}.)\cel{2}Neregulaleso di la movi di la korpo, produkt-}
\l{\cel{3}ata da la perturbeso di la funcioni nervala. - DEFIS.}
\l{}
\l{\sou{ateismo}. Opiniono di la persono qua ne kredas ke existas Deo.}
\l{\cel{3}- DEFIRS.}
\l{}
\l{\sou{ateisto}. La persono qua profesas ateismo. - DEFIRS.}
}
